% =============================================================================
%  CHAPTER 7 -- DISCUSSION                                            [v4 OWNS]
% -----------------------------------------------------------------------------
%  v4.0 (2026-09-24): written by v4 on the v3.98 handover
%  (FROM_v3_v3.98_20260924_201920/README.md): Chapters 5-6 at v3.98, Chapters 1-4 at
%  v2.27 (via v3). The four section labels are the ones Chapters 1 and 4 reference;
%  the subsection labels are new. Every number restates one of Chapter 6, and the
%  \dataref closing each subsection names the tables and sections it comes from.
%  Storyline: ../GUIDE_20260924_results_storyline_author.md (the three steps, the
%  general conclusion); insights: ../INSIGHTS_20260924_results_per_environment.md;
%  vocabulary: ../../Auxiliary/Naming/TRANSLATION_20260914_dev_jargon_to_scientific.md;
%  style: ../../Writing_Hints/HINT_20260920_prompt_is_not_thesis_text.md.
%  Individual changelog: changelogs/v4.0_20260924_init_ch7_ch8_appendix.md
% =============================================================================

\chapter{Discussion}\label{ch:discussion}

The results of \autoref{ch:results} are read here across the five environments: what they say
together (\autoref{sec:disc:interpretation}), where a configuration did not do what it was built for
(\autoref{sec:disc:negative}), what bounds the evidence (\autoref{sec:disc:threats}), and where the
framework itself stops (\autoref{sec:disc:limitations}). Every number in this chapter restates one of
\autoref{ch:results}.

% =============================================================================
\section{Interpretation}\label{sec:disc:interpretation}
% =============================================================================

\subsection{One Result in Three Steps}\label{sec:disc:threesteps}

The study was built in three steps, and the results follow them. The main claim is made and tested on
D3IL-avoiding, the benchmark of \ac{DPCC}, on which the endpoint-projection method was also evaluated.
There the two average-velocity models Pareto-dominate the diffusion model of \ac{DPCC} with its own
temporal U-Net of $4.0$\,M parameters: consistency-interpolated average-velocity matching at one
network evaluation reaches the baseline's success with constraint satisfaction of $1.000$ in $59.2$
control steps against $70.1$, at $18.1$ against $553.4$\,ms per action; analytic average-velocity
matching trades one episode in thirty for the same saving; and no endpoint-projection configuration,
which needs at least three evaluations to guide the sample, improves on that operating point.
D3IL-aligning extends the claim to camera observations and an alignment through contact and reaffirms
it: analytic average-velocity matching dominates the baseline again, nine violation-free contexts
against four, a median box $0.179$ against $0.462$\,m from its target, $266$ against $809$\,ms per
control step, and there endpoint projection is the better projector on the constraint. The quadrotor
scenes are a transfer test: they carry the framework to a new plant and its tracking controller, and
each shows what the plant, the controller or the demonstrations change.

Read together, the five environments sort by their demonstrations. Where the planner learns from
human demonstrations, on D3IL-avoiding and D3IL-aligning, the average-velocity models lead, and the
model and the projection matter together: per-step projection at one evaluation completes
D3IL-avoiding, and on D3IL-aligning, which needs twenty evaluations, endpoint projection keeps every
context violation-free. Where the demonstrations are generated and simple, the three straight lanes
of UAV-corridor and the single route of UAV-s-curve, the objective does not matter: on the corridor the
three flow-based models leave the same violating steps to within $1.6$ per flight at a shared budget
and projection method, and on the s-curve instantaneous-velocity matching leads. UAV-pillars flies the
operating points of D3IL-avoiding on the quadrotor and keeps their result on the declared
constraints; what changes is the plant. And on the quadrotor the plant and its controller matter as
much as the planner: a collision with a pillar the constraint set leaves out ends a flight, and on the
s-curve the same plans succeed or fail by the controller that flies them.
\dataref{\autoref{sec:res:summary:conclusion}; \autoref{tab:avoiding-conclusion};
\autoref{sec:res:aligning:conclusion}; \autoref{sec:res:uav:pillars:conclusion};
\autoref{sec:res:uav:corridor:conclusion}; \autoref{sec:res:uav:scurve:conclusion}.}

\subsection{Where the Saving Comes From}\label{sec:disc:saving}

The step budget moved from training to inference (\autoref{sec:method:fm}), and most of the saving on
D3IL-avoiding is that move. Instantaneous-velocity matching at one evaluation already dominates the
baseline of record, $1.000$ in $67.0$ control steps at $17.3$\,ms, so a flow-based sampler at a small
budget accounts for the thirtieth by itself. The average-velocity objectives add the fewest control
steps of the benchmark, about eight fewer than instantaneous-velocity matching for about a
millisecond more per step, and on D3IL-aligning they add the outcome: at twenty evaluations analytic
average-velocity matching closes $84\,\%$ of the distance to the target, instantaneous-velocity
matching $10\,\%$, and the diffusion baseline's median ends where the box started. The saving does not
rest on the baseline's twenty steps. Retrained at two denoising steps the baseline also reaches
$1.000$, in $61.2$ control steps at $211.4$\,ms, and against that configuration the flow-based models
are twelve times faster rather than thirty.

About half of a projected action's time is the projector, and the price of a projection is set by
what the projector is handed. At one evaluation the flow-based models project a single terminal step,
about $8$ of the $18.3$\,ms analytic average-velocity matching takes. The baseline at two denoising
steps projects a plan one step from noise, at about $45$\,ms per candidate against $2$ to $6$\,ms for
a nearly denoised plan; its two-step time is six times its one-step time. Endpoint
projection shows the same law from the other side: at three evaluations it halves the time per control
step of the two average-velocity models, $74.5$ against $148$\,ms and $76.1$ against $146$\,ms, because
the predicted endpoint it projects is already close to the feasible set. On D3IL-aligning the
activation threshold is compressed as the budget grows, $\eta = 0.2$ at twenty evaluations and $0.4$
at ten, three guiding steps in both cases, and a solve costs about $6$\,ms at twenty evaluations
against $27$\,ms at ten, where the plans handed to the projector are further from their final shape.
\dataref{\autoref{tab:avoiding-dpcc-protocol}; \autoref{tab:avoiding-raw-models};
\autoref{sec:res:avoiding:caveat:noise}; \autoref{tab:avoiding-projectors}; \autoref{tab:va-models};
\autoref{tab:va-threshold}.}

\subsection{The Budget and the Guiding Steps}\label{sec:disc:budget}

Two aims of the thesis pull against each other. Lowering the budget is where the cost saving lives,
and D3IL-avoiding is solved at one evaluation. Endpoint projection needs sampling steps before the
terminal one to guide, $n\sidx{guide} = \lceil \eta\nfe \rceil - 1$ (\autoref{sec:method:degenerate}),
and at one evaluation it has none at any threshold; at two, at the default threshold, none either, and a
run at two evaluations with both projectors given the same candidate plans produced bit-identical
rollouts at $2.45$ times the cost. On D3IL-avoiding endpoint projection therefore has no task to win:
its first admissible budget, three, lies three to eight times to the right of the operating point in
time at the same success. The two aims meet where the task itself needs a larger budget. D3IL-aligning
is run at ten and twenty evaluations because the generative model comes closest to the target there;
the endpoint solve has three guiding steps, and it keeps all ten contexts violation-free where per-step
projection keeps seven to nine, in seven of nine matched comparisons, at a comparable price. Its
advantage is on the constraint, not on the task: it moves the box in about as many contexts, and its
median distances are level with per-step projection's.

On UAV-corridor the two projectors separate only at twenty evaluations, and there they trade. Endpoint
projection leaves fewer and shallower violating steps under both constraints; over the hump with a
single candidate it is also cheaper, $0.2$ against $1.8$ violating steps at $291$ against $355$\,ms, but
its flights are about thirty control steps longer; under the tilt it costs more, $417$ against
$360$\,ms. At three and five evaluations it is as cheap or cheaper and leaves as many violating steps
or more. Neither projector dominates. \guard{On the two average-velocity models the endpoint sampler
queries the network at a zero interval where the plain sampler queries it at the step width
(\autoref{sec:setup:protocol:guiding}), so on those models a difference between endpoint and per-step
projection belongs to the sampler and the point of projection together.}
\dataref{\autoref{sec:res:constraints:degenerate}; \autoref{sec:res:avoiding:projectors};
\autoref{tab:va-projection}; \autoref{tab:va-projection-models}; \autoref{sec:res:uav:projection}.}

\subsection{Candidate Plans and the Selection Rule}\label{sec:disc:selection}

The candidate machinery of \ac{DPCC}, $B = 4$ plans per control step and a rule that picks one, is
where the projector's time goes. Generating four plans instead of one changes the generation time by
at most six per cent, because the four are one batched network evaluation; projecting them
multiplies the projection time by $3.2$ to $4.9$, because the four programs are solved one after
another on the CPU. With a single candidate the three rules select the same plan, and success and
control steps are identical in all fifteen seed--geometry cells of every model. On UAV-corridor the
single-candidate form of endpoint projection at twenty evaluations over the hump costs $291$\,ms per
action against $422$ to $425$\,ms with four candidates, and leaves $0.2$ violating steps against $0.4$
to $0.7$.

Which rule is used matters by model. Cumulative projection cost, the rule \ac{DPCC} selects and the
best one for its diffusion model, is not the flow-based models' rule: on D3IL-avoiding it occasionally
selects a plan that stalls, and analytic average-velocity matching under it takes $72.4$ control steps
at one evaluation and $97.2$ at two, against $58.6$ and $59.4$ under temporal consistency, at $0.933$
against $0.967$ success with constraint satisfaction; the consistency-interpolated model goes from
$60.1$ to $89.7$ steps in the same way. Instantaneous-velocity matching is indifferent, at $1.000$
under all three rules at one and two evaluations. At twenty episodes per seed the rules separate
further: the operating point reads $298$ of $300$ under temporal consistency against $279$ and $283$
under the other two. On UAV-s-curve the three rules cross the finish line on five flights of ten each.
Where the rule selects the same plan, as with one candidate, or where the model is indifferent to it,
the machinery can be dropped and the projection time with it; where the flow-based models are operated
with four candidates, temporal consistency is their rule, and solving the four programs in parallel
would take analytic average-velocity matching at one evaluation from $18.1$ to about $11.7$\,ms per
control step, a prediction from two measured points.
\dataref{\autoref{sec:res:ablations}; \autoref{tab:avoiding-dpcc-protocol};
\autoref{app:avoiding-twenty}; \autoref{app:uav-corridor-rules}; \autoref{app:uav-scurve-rules}.}

\subsection{What the Demonstrations Decide}\label{sec:disc:demonstrations}

In every environment the constraints are absent from the demonstrations, and satisfying them is the
projector's work. On D3IL-avoiding no unprojected row of any model exceeds $0.133$ success with
constraint satisfaction, and $15.5$ to $19.0$ control steps of an average episode violate a constraint
at every budget; under per-step projection the same models read $0.967$ to $1.000$. On UAV-corridor a
flight violates the tilt on $74.5$ to $99.4$ control steps before projection and on $1.4$ to $10.4$
after it, and the projected flights descend under the tilt on $540$ of $560$ and climb over the hump on
all $560$ where the unprojected ones fly level through both. The concept of \ac{DPCC}, a learned
planner whose plans a projector holds inside constraints the data never satisfied, carries over from
the table to a three-dimensional scene and to constraints that leave the horizontal plane.

What the demonstrations decide is what the generative model can show. The human demonstrations of the
two D3IL tasks are multimodal, and on them the objectives separate: on D3IL-aligning analytic
average-velocity matching ends closer to the target than the baseline in all ten contexts and than
instantaneous-velocity matching in nine. The generated demonstrations of the quadrotor scenes carry
one path forward from any position, and on them the objectives do not separate. On the corridor the
four models cannot be told apart before projection, and after it the three flow-based objectives agree
on violating steps to within $1.6$ per flight; the budget and the projection method order the
configurations. On the s-curve, one route of straight segments joined by arcs, instantaneous-velocity
matching crosses the finish line most often, nine flights of ten at one evaluation against six for the
consistency-interpolated model and three for the analytic one. The demonstrations also bound the plan.
On the corridor the plans never rise above $1.30$\,m, the top of the demonstrated altitude range, while
the apex of the hump needs $1.48$\,m, and per-step projection at one and two evaluations stops $177$ of
$180$ flights short of it. On the s-curve every commanded path cuts the second inside corner of the
crossover by $8$ to $19$\,cm and crosses the gap about $0.2$\,m later than the demonstrated route. And on
UAV-pillars the plans pass the pillars as closely as the human demonstrations passed the posts, at the
rod's radius, which the scale of the scene makes the rotor reach.

The corridor's remaining violations are structural. Every violating step of every projected flight
under the tilt lies in the plane's last $0.3$\,m, and at the first of them the commanded position is
already past the plane's end, $0.49$ to $0.61$\,m ahead of the vehicle and clear of the constraint. The
span of a plane is checked at the commanded position, so the plan is released before the vehicle has
left the plane; the budget makes the residue shallower, from $8.8$ to $11.9$\,cm at one and two
evaluations to $1.0$ to $2.2$\,cm at twenty, and does not remove it.
\dataref{\autoref{tab:avoiding-raw-models}; \autoref{tab:avoiding-dpcc-protocol};
\autoref{tab:uav-corridor-raw}; \autoref{tab:uav-corridor}; \autoref{sec:res:uav:projected};
\autoref{tab:va-models}; \autoref{tab:uav-scurve}; \autoref{sec:res:uav:pillars};
\autoref{tab:uav-demos}.}

\subsection{The Plant and Its Controller}\label{sec:disc:plant}

On the manipulator a plan reaches the joints through inverse kinematics and a joint-space PD law, and
the end effector holds its setpoint; on the quadrotor it reaches the rotors only through the cascaded
geometric controller, and the vehicle holds nothing by itself. UAV-pillars flies the Pareto-optimal
operating points of D3IL-avoiding with that plant exchanged and nothing else. The declared constraints
transfer exactly: after projection no flight of the $78$ that reach the goal line violates a halfspace
or enters the keep-out disk, the control steps agree to within three and the time per action to within
a millisecond, and the two average-velocity models keep their order and their gap, the
consistency-interpolated one a flight in thirty ahead at both budgets. The vehicle tracks more closely
than the arm, within about $0.007$ of the commanded position in the table's units against the
manipulator's $0.02$ to $0.04$, well inside the tightening margin of $0.025$. What the plant changes is
the cost of an obstacle. Five of the six pillars are in no declared constraint set, the plans skim
them, and a quadrotor whose rotor strikes a pillar loses attitude control: $42$ of the $120$ projected
flights end that way, on $50$ of the $51$ collisions with the commanded path itself inside the rotor
reach, and success within the constraints falls from $0.967$--$1.000$ to $0.600$--$0.700$. Both models
lose the same share on the same geometry, so the collisions do not re-order them. A declared constraint
set transfers across the plant; freedom from collision with what it leaves out does not.

UAV-s-curve shows the controller. The same plans of instantaneous-velocity matching at one evaluation
cross the finish line on nine flights of ten under either controller without projection, and under
per-step projection on five of ten under the cascaded geometric controller and on none under MuJoCo
MPC, seven of whose flights invert against the outer wall of the second straight. The unprojected plans
violate on $23.0$ control steps per flight under the cascaded law and on $41.5$ under MuJoCo MPC, which
follows the commanded path less closely, $5$\,cm on average and up to $29$\,cm against about $1$\,cm. On
the configuration of record the cascaded geometric controller is the better one, and it is twenty-four
times cheaper: the controller and the simulator step cost $5.2$\,ms per executed control step against
$125.4$\,ms for MuJoCo MPC, which solves an optimisation problem at every control step where the other
evaluates an algebraic law. MuJoCo MPC's one gain, keeping every unprojected flight upright and within
$0.3$\,m of the goal point, does not survive projection. The budget adds a caveat of its own on this
scene: for every flow-based model the crossings fall as the budget grows, nine to seven to six for
instantaneous-velocity matching, six to one for the consistency-interpolated model, three to none for
the analytic one, while the plans grow smoother; the rough plans of one evaluation, re-drawn at every
control step, vary more from step to step, and on this route that variation appears to help the
controller through the turns.
\dataref{\autoref{tab:uav-pillars-raw}; \autoref{tab:uav-pillars-geo};
\autoref{sec:res:uav:pillars:conclusion}; \autoref{tab:uav-controller};
\autoref{tab:uav-controller-cost}; \autoref{tab:uav-scurve}; \autoref{fig:uav-scurve-plans}.}

% =============================================================================
\section{Negative and Inconclusive Results}\label{sec:disc:negative}
% =============================================================================

\paragraph{The consistency-interpolated objective does not improve with the budget.} On D3IL-aligning
its median final distance moves from $0.3738$\,m at two evaluations to $0.3897$\,m at twenty, against
the analytic objective's $0.2779$ to $0.0741$\,m, and it reaches $60\,\%$ of the distance only at a
hundred evaluations; the more aggressive setting $\alpha\sidx{end} = 0.05$ ends further away still,
$0.4488$\,m at twenty. Its training target is built from a stop-gradient copy of the network, so part
of the network's own error enters the field it learns, and smaller steps do not remove it. A task that
rewards more evaluations exposes this; a task saturated at one evaluation does not, and on
D3IL-avoiding the same objective is the model with the highest success with constraint satisfaction,
$1.000$ against the analytic objective's $0.967$, one episode in thirty. It is ahead of the analytic
objective by one flight in thirty on UAV-pillars and by three crossings on UAV-s-curve, level with it
on UAV-corridor, and behind it by most of the task on D3IL-aligning. The consistency interpolation is
therefore not a general improvement on the analytic target, and the analytic target is the one that
leads or holds in every environment.

\paragraph{Instantaneous-velocity matching and the diffusion baseline are level on outcome.} Where
the task separates the objectives, on D3IL-aligning, the two have no ordering: their medians are
$0.4085$ and $0.4619$\,m, each ends closer in about half of the ten contexts, and they leave the box
unmoved in four and five of them. On D3IL-avoiding both reach $1.000$ under the projector, and on
UAV-corridor before projection their violating steps at twenty evaluations lie within one standard
deviation of each other, $94.7$ against $99.4$ under the tilt and $31.3$ against $28.6$ over the hump;
after projection the baseline keeps fewer, $0.8$ and $0.0$ against $2.1$ and $1.8$, at $666$ and
$654$\,ms per action against $360$ and $355$. What separates the two models is cost: a thirtieth of the
time per action on D3IL-avoiding at one evaluation, a ninth to a twenty-fourth on the corridor.
\guard{Instantaneous-velocity matching is sampled from initial noise of half its training scale in
every environment (\autoref{sec:setup:protocol:avoiding}); its order against the average-velocity
models on D3IL-aligning and the quadrotor scenes carries that difference.}

\paragraph{Endpoint projection has nothing to win at the small budgets.} On D3IL-avoiding the
operating point is one evaluation, where the method has no guiding step, and its first admissible
budget, three, is off the cost frontier. On UAV-s-curve the selected configuration is at one evaluation
and endpoint projection is not flown; on UAV-pillars the operating points are at one and two. Where the
budget is large it is the better projector on the constraint on D3IL-aligning and a trade-off on
UAV-corridor. What endpoint projection contributes is therefore conditional on the budget the task
needs, and D3IL-avoiding, the benchmark of \ac{DPCC}, is not a task that needs one.

\paragraph{No context of D3IL-aligning ends in position.} The supported combination keeps all ten
contexts violation-free, and the box ends within the $1.8$\,cm of the task's own position test in none
of them: $2.2$ and $4.7$\,cm away in the best two, $13$ to $27$\,cm in six, unmoved in two. The
unprojected plan reaches $2.8$ and $2.9$\,cm in its best two. Pursuing no violation costs completion on
this task: the projected median ends at $0.197$\,m against $0.090$\,m unprojected.

\paragraph{No flight of UAV-corridor under the tilt and none of UAV-s-curve is violation-free.} Under
the tilt every projected flow-based flight cuts into the plane over its last $0.3$\,m, the hand-over
residue of \autoref{sec:disc:demonstrations}; only the baseline's flights are violation-free there,
three to eight of ten, at twenty-four times the time per action. On the s-curve every commanded path of
every configuration cuts the second inside corner, and per-step projection, which binds the measured
position and not the commanded one, leaves the cut in place, $8$ to $9$\,cm deep, while costing four of
the nine crossings and $126$ to $155$\,ms per control step. MuJoCo MPC removes the one inversion of the
unprojected plans, adds violations, and loses every projected flight.

\paragraph{A third of the UAV-pillars flights collide.} With the declared constraints satisfied on
every flight, $42$ of $120$ projected flights end in a collision with a pillar outside the constraint
sets, concentrated on \emph{both-hard}, whose two halfspaces funnel the plan past the outer pillars of
the third row: two to four of ten flights succeed there against seven to nine on the single-halfspace
geometries.

\paragraph{More evaluations are not always better.} On D3IL-aligning a hundred evaluations reduce the
analytic objective's median by $0.0068$\,m at $4.7$ times the cost of twenty. On D3IL-avoiding
instantaneous-velocity matching at twenty evaluations holds $1.000$ at $476.7$\,ms per action, about
$28$ times its cost at one evaluation, for four fewer control steps. On UAV-s-curve the crossings fall
with the budget for every flow-based model, and the flights that are lost invert at the second turn.
\dataref{\autoref{tab:va-models}; \autoref{sec:res:aligning:consistency};
\autoref{tab:avoiding-dpcc-protocol}; \autoref{tab:uav-corridor-raw}; \autoref{tab:uav-corridor};
\autoref{sec:res:summary:conclusion}; \autoref{tab:uav-scurve-projection};
\autoref{tab:uav-controller}; \autoref{tab:uav-pillars-geo}; \autoref{tab:uav-scurve}.}

% =============================================================================
\section{Threats to Validity}\label{sec:disc:threats}
% =============================================================================

\paragraph{Seeds and sample sizes.} D3IL-avoiding and UAV-pillars are evaluated on five training
seeds; D3IL-aligning, UAV-corridor and UAV-s-curve on one (\autoref{tab:protocol}). The variation
across seeds is shown on D3IL-avoiding: the flow-based models' spread over the seeds stays within
$0.075$ in success with constraint satisfaction and within about three control steps, the baseline's
reaches $0.391$ at one denoising step and $8.2$ control steps at its own budget. At two episodes per
seed and geometry a rate moves in steps of one episode in thirty; the same configurations at twenty
episodes per seed read within that step, $298$ of $300$ for the operating point against $29$ of $30$,
and the standard deviation over the seeds falls from $0.075$ to $0.015$ (\autoref{app:avoiding-twenty}).
Ten contexts and ten flights resolve to one in ten, and a one-context difference on D3IL-aligning is
within what a re-run changes: a second evaluation of one projected cell moved it by one context and
$3$\,cm in median.

\paragraph{The alignment contexts.} The ten contexts of D3IL-aligning are drawn from the training
contexts of the D3IL task, so its results measure how well each model plans for initial poses of the
kind it was trained on. On the tightened set one of the ten is never attempted: the tightening enlarges
the keep-out region until it overlaps the box of context~8, and every $x/10$ of the tightened alignment
tables is $x-1$ of nine attempted contexts. Comparisons between rows are unaffected.

\paragraph{What the four models do not share.} The models share the backbone and its parameter count,
the data, the plan horizon and the projection they are read through. They do not share their
optimiser settings, which were set per objective (\autoref{tab:train}), and they do not share their
sampler's prior. Instantaneous-velocity matching is sampled from noise of standard deviation $0.5$,
the scale inherited from the baseline's sampler, while it was trained from noise of standard deviation
$1$; the two average-velocity models are sampled at the scale they were trained at. Wherever a
difference between instantaneous-velocity matching and an average-velocity model is read as a property
of the objective, the prior scale is a second difference between them. On D3IL-avoiding
instantaneous-velocity matching holds $1.000$ under every rule at one and two evaluations, so the
ordering there does not rest on it. Where endpoint projection is compared with per-step projection on
the two average-velocity models, the endpoint sampler queries the network at a zero interval where the
plain sampler queries the step width, so that difference belongs to the sampler and the point of
projection together.

\paragraph{What the time per action contains.} The time to compute one action is measured around the
planner: generating the candidate plans, projecting them and selecting one
(\autoref{sec:setup:metrics:avoiding}). It contains neither the inverse kinematics and the joint-space
controller of the manipulator, nor the geometric controller of the quadrotor, nor the simulator step,
and the simulation waits for the planner, so no plan interval of this thesis is a rate the loop meets in
real time. The one place the execution stack is priced is the controller comparison of UAV-s-curve: on
the selected plans without projection an executed control step costs $14.0$\,ms under the cascaded
geometric controller, of which the planner is $8.9$\,ms and the controller with the simulator step
$5.2$\,ms, and $135.8$\,ms under MuJoCo MPC, of which $125.4$\,ms are the controller. At the
one-evaluation operating point the cascaded controller's share is of the planner's order; against the
diffusion baseline's $553.4$\,ms per action on D3IL-avoiding it is a hundredth. The manipulator's
inverse kinematics and joint controller were not timed. All times were measured on one shared node,
whose GPU a job holds alone while its CPU, memory and I/O are shared, and the MuJoCo rollout and the
SLSQP projection run on the CPU (\autoref{app:repro}); they are compared between configurations
measured on that node and are not hardware benchmarks.

\paragraph{The projection stacks of the quadrotor scenes differ.} UAV-corridor constrains the
commanded position in addition to the measured one and leaves the action unconstrained; UAV-s-curve
binds the measured position and the action, as D3IL-avoiding does (\autoref{sec:setup:protocol:uav}).
The two scenes are not pooled, and the s-curve's corner cut is read under its own stack: a projector
that also binds the commanded position was not flown there.

\paragraph{The quadrotor demonstrations are generated.} The demonstrations of UAV-corridor and
UAV-s-curve are reference paths laid out analytically and flown in MuJoCo by the same controller the
evaluation uses; expert means a competent generated flight and nothing more
(\autoref{sec:setup:data:uav}). The corridor's three lanes and the s-curve's single route are what the
planners learn, and the tie among the objectives on the corridor and the order on the s-curve are read
on those data. UAV-pillars has no demonstrations of its own and flies the planner of D3IL-avoiding.

\paragraph{The controller comparison.} MuJoCo MPC replaces the cascaded geometric controller on one
configuration, ten flights against ten, with both controllers at their nominal settings; the gains of
the cascaded controller were not tuned, and the two alternative settings that exist scale its position
loop by $\pm20\,\%$ (\autoref{sec:res:uav:controller}). What the comparison supports is which of the
two, as configured, flies the scene better, and not that no tuning of either could do more.
\dataref{\autoref{tab:protocol}; \autoref{sec:setup:metrics:spread}; \autoref{tab:avoiding-dpcc-protocol};
\autoref{app:avoiding-twenty}; \autoref{sec:res:aligning:projection:models};
\autoref{sec:setup:protocol:avoiding}; \autoref{sec:setup:protocol:guiding};
\autoref{tab:uav-controller-cost}; \autoref{sec:setup:protocol:uav}; \autoref{sec:setup:data:uav}.}

% =============================================================================
\section{Limitations}\label{sec:disc:limitations}
% =============================================================================

\paragraph{The plan is a sequence of positions.} The generative model plans position increments and
never an actuator command; everything below the increment is the plant's controller
(\autoref{sec:method:setpoint}). The projection's dynamics model is a first-order integrator over the
commanded quantity, so feasibility is a statement about that model and its mismatch bound, which the
tightening absorbs, and not about the plant. On the manipulator this holds, because the end effector
tracks its setpoint within the margin; on the quadrotor whether a plan can be flown is the
controller's, which is what UAV-s-curve shows.

\paragraph{A constraint is what is declared.} The projector holds a plan inside the declared geometry,
halfspaces and keep-out regions, and protects nothing else. On UAV-pillars five of the six pillars are
in no constraint set and a third of the flights end on them. Whether a scene is safe follows from its
constraint set and not from the projection; the remedy is a constraint per obstacle, or a margin beyond
the rotor reach, and it is a decision of scene design. The span over which a wall or a plane acts is
checked at the commanded position, half a metre ahead of the vehicle on the corridor, which leaves the
hand-over residue at the end of every plane.

\paragraph{Endpoint projection is defined on a velocity field.} It adds its control to the field of a
flow-based sampler and has no counterpart for the diffusion model, whose reverse chain has none
(\autoref{sec:method:hardflow}). The comparison of the two projection methods is therefore made within
the flow-based family, and the diffusion baseline is read under per-step projection alone.

\paragraph{Plan quality is read through violating steps and the drawn plans.} The plans are scored by
success, constraint satisfaction, control steps and time, and inspected as drawn (\autoref{fig:raw-plans},
\autoref{fig:uav-scurve-plans}). No curvature, jerk or path-length measure is taken over the stored
plans, so the ribbon of the average-velocity models against the scribble of the one-step baseline, and
the rough low-budget plans that cross the s-curve most often, are read from the drawings and the counts
beside them.

\paragraph{Scope of the transfer.} The vision-conditioned task adds camera observations and contact to
the manipulator in one step; the state-observed form of the alignment task was not run, so what
D3IL-aligning shows is that the result survives a harder, differently observed task, and not the step
from state to vision alone. The quadrotor is observed through its state, at one altitude on UAV-pillars,
where the plan carries no altitude, and within the $0.90$ to $1.30$\,m of the demonstrations elsewhere.
No loop of this thesis was run in real time, and no environment is a physical system.
\dataref{\autoref{sec:method:manipulator}; \autoref{sec:method:env:uav}; \autoref{sec:res:uav:pillars};
\autoref{sec:res:uav:projected}; \autoref{sec:method:hardflow}; \autoref{tab:uav-demos};
\autoref{sec:method:setpoint}.}
